\documentclass{situationist-preprint}

\title{The Creation of Fictional Characters}
\msid{LNG-14-b}
\msfamily{LNG-14}
\msdiscipline{Philosophy of Language}
\mstype{Companion manuscript}
\mskeywords{fiction; fictional characters; abstract artefacts; creation; pretense; dependence}

\begin{document}
\maketitle

\begin{abstract}
LNG-14-a argues that fictional names refer to abstract artefacts created by authors. We say what the creating consists in. A character is created when an author, composing a work of fiction, makes pretended reference with a name or description and so makes it the case that the work ascribes properties to something. Creation does not require publication, completion or any further act; it depends historically on the author's act and is individuated by it, so that two authors who independently write the same story create two characters. The account removes the appearance that creation by writing is mysterious and fixes the dates of characters' existence.
\end{abstract}

\section{What the author does}

On the view of LNG-14-a, \enquote{Holmes} refers to an abstract artefact, an object brought into existence by the intentional activity of an agent and not located in space (Definition 2.2 of LNG-14-a). Conan Doyle did not assert that there was a detective called Holmes; he pretended to refer to one and pretended to assert things about him \citep{Searle1975}. We propose that this pretense is the creating act.

\begin{principle}[Creation]\label{p:create}
If an author, in composing a work of fiction, makes pretended reference with a name or description, and no existing character is thereby picked out, then a character comes into existence that the work ascribes the pretended properties to; and every fictional character comes into existence in this way.
\end{principle}

The principle states sufficient and necessary conditions in terms that involve no special metaphysical power. The author does something ordinary, writing sentences in a pretending mode, and the existence of the character is a consequence of that act given the practice of fiction, as the existence of a contract is a consequence of signatures given the practice of law, or the existence of a new chess opening of a sequence of moves given the practice of chess. \citet{Thomasson1999} calls such objects dependent abstracta: they depend on intentional acts for their coming into existence and on something concrete, copies of the work or the memories of readers, for their continuing to exist. This answers the objection from the mystery of creation that LNG-14-a addresses: the mystery is no greater than that of any social artefact.

\section{Consequences}

\paragraph{Dates.} Holmes came into existence when Conan Doyle first made pretended reference with \enquote{Holmes} in composing \emph{A Study in Scarlet}, not at its publication. Characters of unpublished drafts exist; characters of stories never written down, if their author composed them in thought with the same pretense, exist too. What does not exist is a character of a story no one has composed.

\paragraph{Individuation.} Since creation is an act, characters are individuated by their creating acts. Two authors who, without contact, compose word-for-word identical stories create two characters, which share all the properties their works ascribe and differ in their authors. Conversely, an author who continues a character in later works does not create a new one: the pretended reference in the later works is to the existing character, as the proviso of Principle~\ref{p:create} states. This is why the later Holmes stories ascribe their further properties to the same object to which, by Assumption 2.8 of LNG-14-a, the first ascribes detection.

\paragraph{Ascription grows; creation does not repeat.} Holmes is ascribed different properties as the canon grows. The character does not change in its intrinsic nature, having none to change; the relation of ascription to it holds between more works and more properties. A character thus exists from its first pretended reference, and what is true of it according to the stories is fixed by all the works that refer to it.

\begin{objection}
Pretended reference is not reference. If Conan Doyle failed to refer when he wrote \enquote{Holmes}, his act cannot have brought into being something to which \enquote{Holmes} refers.
\end{objection}

\begin{reply}
The act need not refer to what it creates, any more than a signature refers to the contract it brings about. Once the character exists, critics and readers who use \enquote{Holmes} in critical discourse refer to it, and their reference is secured by a chain leading back to the creating act, in the manner of \citet{Kripke2013}, who also holds that fictional characters are abstract entities created by works of fiction. The author's own uses are pretense; the uses that refer are those that follow.
\end{reply}

\section{Conclusion}

A fictional character is created by an author's pretended reference in composing a work. This identifies the creating act, explains why creation is no more mysterious than the creation of other social artefacts, and settles when characters begin to exist and how they are counted. It completes the account of LNG-14-a, on which talk about Holmes is talk about an existing abstract artefact.

\begin{thebibliography}{9}\small
\bibitem[Kripke(2013)]{Kripke2013} Kripke, S. (2013). \emph{Reference and Existence: The John Locke Lectures}. Oxford University Press.
\bibitem[Searle(1975)]{Searle1975} Searle, J.~R. (1975). The logical status of fictional discourse. \emph{New Literary History}, 6(2), 319--332.
\bibitem[Thomasson(1999)]{Thomasson1999} Thomasson, A.~L. (1999). \emph{Fiction and Metaphysics}. Cambridge University Press.
\end{thebibliography}

\end{document}
